\section{答题时间规划与答题技巧}

\subsection{试卷结构速览}

\begin{center}
\begin{tabular}{c|c|c|c}
\textbf{题型} & \textbf{题号} & \textbf{分值} & \textbf{建议用时} \\ \hline
现代文阅读 I（信息类） & 1--5 & 19 分 & 25 分钟 \\
现代文阅读 II（文学类） & 6--9 & 16 分 & 20 分钟 \\
文言文阅读 & 10--14 & 20 分 & 20 分钟 \\
古代诗歌鉴赏 & 15--16 & 9 分 & 10 分钟 \\
名篇名句默写 & 17 & 6 分 & 3 分钟 \\
语言文字运用 & 18--22 & 20 分 & 17 分钟 \\
写作 & 23 & 60 分 & 55 分钟 \\ \hline
\textbf{总计} & \textbf{23 题} & \textbf{150 分} & \textbf{150 分钟}
\end{tabular}
\end{center}

\subsection{答题顺序策略}

\textbf{总体原则}：先易后难，遇障则绕。作文务必留足 50--55 分钟。

\textbf{推荐顺序}：

\begin{enumerate}
  \item \textbf{名篇名句默写（17 题，3 分钟）} —— 开卷即抢分，趁头脑清醒默写
  \item \textbf{语言文字运用（18--22 题，17 分钟）} —— 题型固定，难度适中
  \item \textbf{现代文阅读 I 信息类（1--5 题，25 分钟）} —— 速读定位，选择为主
  \item \textbf{文言文阅读（10--14 题，20 分钟）} —— 看题读文，先断后译
  \item \textbf{古代诗歌鉴赏（15--16 题，10 分钟）} —— 读诗析题，按模板作答
  \item \textbf{现代文阅读 II 文学类（6--9 题，20 分钟）} —— 细读赏析，规范表述
  \item \textbf{写作（23 题，55 分钟）} —— 审题立意 5--8 分钟 + 提纲 5 分钟 + 成文 35 分钟 + 检查 5 分钟
  \item \textbf{剩余时间检查} —— 重点检查默写（易错字）、选择题涂卡
\end{enumerate}

\subsection{各题型时间分配细则}

\subsubsection{论述类／信息类阅读（25 分钟）}

\begin{itemize}
  \item 浏览全文：3 分钟（抓中心句、论点论据）
  \item 选择题（1--3 题）：10 分钟（逐项比对，注意偷换概念、以偏概全等陷阱）
  \item 简答题（4--5 题）：12 分钟（概括内容、分析论证特点）
\end{itemize}

\subsubsection{文学类阅读（20 分钟）}

\begin{itemize}
  \item 通读全文：5 分钟（把握情节、人物、主题）
  \item 选择题（6--7 题）：5 分钟（艺术特色判断）
  \item 简答题（8--9 题）：10 分钟（作用分析、含义理解、手法赏析）
\end{itemize}

\subsubsection{文言文阅读（20 分钟）}

\begin{itemize}
  \item 断句题（10 题）：3 分钟（通语感、找标志词）
  \item 词语解说（11 题）：3 分钟（结合语境排除）
  \item 内容分析（12 题）：4 分钟（比对原文细节）
  \item 翻译（13 题）：6 分钟（字字落实 + 语意通顺）
  \item 简答（14 题）：4 分钟（筛选概括）
\end{itemize}

\subsubsection{诗歌鉴赏（10 分钟）}

\begin{itemize}
  \item 读诗 + 选择题（15 题）：4 分钟
  \item 简答题（16 题）：6 分钟（手法 + 内容 + 情感三步骤）
\end{itemize}

\subsubsection{写作（55 分钟）}

\begin{itemize}
  \item 审题立意：5--8 分钟（圈画关键词，确立中心论点）
  \item 列提纲：5 分钟（分论点 + 论据素材）
  \item 正文写作：35 分钟（开头 5 分钟 + 主体 25 分钟 + 结尾 5 分钟）
  \item 检查修改：5 分钟（错别字、语病、字数达标）
\end{itemize}

\subsection{考场应急策略}

\begin{itemize}
  \item \textbf{时间不足}：优先保作文（至少 45 分钟）；选择题快速排除法；简答题列要点而非展开
  \item \textbf{文言文卡壳}：先看第 12 题内容分析题，反向理解文章大意
  \item \textbf{作文跑题风险}：反复读材料，圈关键词；开头点明中心论点；每段段首设分论点句
  \item \textbf{默写遗忘}：跳过，做完其他题后再回忆；可用同篇目其他句子替代
  \item \textbf{选择题犹豫}：第一直觉往往准确，除非确知错误否则不改
\end{itemize}
